\documentclass[12pt]{article}

\usepackage[utf8]{inputenc}
\usepackage{latexsym,amsfonts,amssymb,amsthm,amsmath}
\usepackage{tikz}
\usepackage{stanli}
\usetikzlibrary{shapes,calc,through,3d,patterns}
\usetikzlibrary{arrows.meta,decorations.pathreplacing,angles,quotes}
\usepackage{multicol}
\usepackage{geometry}
\usepackage{mathtools}
\usepackage{siunitx}
\usepackage{booktabs}

\setlength{\parindent}{0in}
\setlength{\oddsidemargin}{-0.25in}
\setlength{\textwidth}{7in}
\setlength{\textheight}{9.3in}
\setlength{\topmargin}{-1in}
\setlength{\headheight}{18pt}

\newcommand\blfootnote[1]{%
    \begingroup
    \renewcommand\thefootnote{}\footnote{#1}%
    \addtocounter{footnote}{-1}%
    \endgroup
}

\definecolor{darkred}{rgb}{0.6,0,0}

\newcommand{\bolddelta}{\boldsymbol\delta}
\newcommand{\boldxi}{\boldsymbol\xi}

\newcommand{\T}{^{T}}
\newcommand{\bK}{\overline{\mathbf K}}
\newcommand{\bM}{\overline{\mathbf M}}
\newcommand{\vct}[1]{\mathbf{#1}}
\newcommand{\colv}[1]{\begin{bmatrix}#1\end{bmatrix}}

\begin{document}

\noindent ME473 - Dynamic finite element analysis of structures\hfill EPFL\\
Stefano Burzio \hfill 2025\\
\vspace{0.3cm}

\hspace*{\fill}\textbf{\Large Problem set 10 - solutions}\hspace*{\fill}

\hrulefill

\subsection*{Problem 1}
\begin{itemize}
   \item[a)] Write $c_\theta=\cos\theta$, $s_\theta=\sin\theta$. The rotated off-diagonal entry is the stiffness product of the first
         and second columns of $\vct R$:
         \begin{align*}
            k'_{ij}
             & = (c_\theta,s_\theta)\vct K_{(i,j)}(-s_\theta,c_\theta)\T  \\
             & =k_{ij}(c_\theta^2-s_\theta^2)
            -(k_{ii}-k_{jj})c_\theta s_\theta                             \\
             & =k_{ij}\cos(2\theta)-\tfrac12(k_{ii}-k_{jj})\sin(2\theta).
         \end{align*}
         Setting $k'_{ij}=0$ gives, for a nonzero pivot and unequal diagonal entries,
         \[
            \tan(2\theta)=\frac{2k_{ij}}{k_{ii}-k_{jj}}.
         \]
         If $k_{ij}=0$, take $\theta=0$. If $k_{ii}=k_{jj}$ and $k_{ij}\ne0$,
         take $\theta=\pi/4$; then $\cos(2\theta)=0$. In all cases
         $|\theta|\leq\pi/4$. For the mass block, $\vct M_{(i,j)}=\vct I$, we obtain
         \[
            \vct M_{(i,j)}'=\vct J\T\vct M_{(i,j)}\vct J
            =\begin{bmatrix}c_\theta&s_\theta\\-s_\theta&c_\theta\end{bmatrix}
            \begin{bmatrix}1&0\\0&1\end{bmatrix}
            \begin{bmatrix}c_\theta&-s_\theta\\s_\theta&c_\theta\end{bmatrix}
            =\begin{bmatrix}c_\theta^2+s_\theta^2&0\\0&c_\theta^2+s_\theta^2\end{bmatrix}
            =\vct I,
         \]
   \item[b)] The two columns of the rotation matrix $\vct J(\alpha,\gamma)$ in the generalized Jacobi method are $\vct u=(1,\gamma)\T$ and $\vct v=(\alpha,1)\T$. For any symmetric block $\vct A_p$, we have
         \[
            \vct u\T\vct A_p\vct v
            = a_{ii}\alpha+a_{jj}\gamma+a_{ij}(1+\alpha\gamma).
         \]
         Applying this product to $\vct K_p$ and $\vct M_p$ and setting each result to zero gives the two equations in the question.
   \item[c)] Denote the two pivot equations by $F_K=0$ and $F_M=0$, in the order given:
         \begin{align*}
            F_K=k_{ii}\alpha+k_{jj}\gamma+k_{ij}(1+\alpha\gamma) & =0, \\
            F_M=m_{ii}\alpha+m_{jj}\gamma+m_{ij}(1+\alpha\gamma) & =0.
         \end{align*}
         The combination $m_{ij}F_K-k_{ij}F_M$ gives
         $\bar k_{ii}\alpha+\bar k_{jj}\gamma=0$. Two other combinations,
         $m_{jj}F_K-k_{jj}F_M$ and $k_{ii}F_M-m_{ii}F_K$, give
         \[
            2\bar k\,\alpha=\bar k_{jj}(1+\alpha\gamma),\qquad
            2\bar k\,\gamma=-\bar k_{ii}(1+\alpha\gamma).
         \]
         Write $\alpha=\bar k_{jj}/x$ and $\gamma=-\bar k_{ii}/x$. Since
         $1+\alpha\gamma=1-\bar k_{ii}\bar k_{jj}/x^2$, substitution gives
         \[
            x^2-2\bar k\,x-\bar k_{ii}\bar k_{jj}=0,\qquad
            x_\pm=\bar k\pm\sqrt{\bar k^2+\bar k_{ii}\bar k_{jj}} .
         \]
         The larger-magnitude root adds terms of the same sign, avoiding
         subtractive cancellation. For a nondegenerate pivot this gives
         \[
            x=\bar k+\operatorname{sgn}(\bar k)
            \sqrt{\bar k^2+\bar k_{ii}\bar k_{jj}},\qquad
            \alpha=\frac{\bar k_{jj}}{x},\qquad
            \gamma=-\frac{\bar k_{ii}}{x}.
         \]
         When $\bar k=0$ and the radicand is positive, the roots have equal
         magnitude. Take the positive root, equivalently $\operatorname{sgn}(0)=+1$
         for this root-selection rule. An already diagonal pair needs no update.
         If all three barred quantities vanish for a coupled pivot, the blocks are
         proportional. In this case take $\gamma=0$ and
         $\alpha=-m_{ij}/m_{ii}$. This diagonalizes $\vct M_p$ and its multiple
         $\vct K_p$ without dividing by $x=0$.
   \item[d)] Let $\vct M_p=\vct I$ to recover the standard rotation matrix $\vct J(\theta)$ notice that the mass equation $F_M=0$ gives $\alpha=-\gamma$ Putting $t=\gamma$ simplify the stiffness equation $F_K=0$ to
         \[
            (k_{ii}-k_{jj})t=k_{ij}(1-t^2).
         \]
         The columns have the same length $\sqrt{1+t^2}$. Normalizing both gives
         \[
            \frac{1}{\sqrt{1+t^2}}\colv{1&-t\\t&1}
            =\colv{\cos\theta&-\sin\theta\\\sin\theta&\cos\theta},
            \qquad t=\tan\theta .
         \]
         For $|t|<1$, the identity $\tan(2\theta)=2t/(1-t^2)$ gives the formula in part (a). With equal diagonals and a nonzero pivot, $t=\pm1$ gives the endpoint rotations $\theta=\pm\pi/4$.
\end{itemize}


\subsection*{Problem 2}

\textbf{1. Formulation of the eigenvalue problem.} We consider the longitudinal vibrations of a uniform free-free bar. The bar is discretized using a single quadratic finite element with three equally spaced nodes. Its dynamic behavior is described by the generalized eigenvalue problem:
\[
   \mathbf{K} \mathbf{p} = \lambda \mathbf{M} \mathbf{p},
\]
where \( \lambda = \omega^2 \) are the squared natural frequencies and \( \mathbf{p} \) are the corresponding mode shapes. The stiffness and lumped mass matrices are defined as:
\[
   \mathbf{K} = \frac{EA}{3\ell}
   \begin{bmatrix}
      7  & -8 & 1  \\
      -8 & 16 & -8 \\
      1  & -8 & 7
   \end{bmatrix}, \quad
   \mathbf{M} = \frac{\rho A \ell}{6} \cdot \text{diag}(1, 4, 1).
\]

\begin{verbatim}
    % Parameters (symbolic for generality, assign numerical values as needed)
    clear
    syms E A rho l real

    % Exact eigenvalues 
    Lambda_exact = [2*E / (rho * l^2); 14*E / (rho * l^2)];

    % Exact eigenvectors (modes)
    p1 = [-1; -1; -1];
    p2 = [-1; 0; 1];

    P_exact = [p1, p2];

    % Stiffness matrix K
    K = (E*A)/(3*l) * [  7  -8   1;
            -8  16  -8;
            1  -8   7 ];

    % Lumped mass matrix M
    M = (rho*A*l)/6 * diag([1, 4, 1]);
\end{verbatim}

\textbf{2. Spectral shift for rigid-body mode.} To improve convergence, we apply a spectral shift to eliminate the rigid-body mode:
\[
   \mathbf{K}_\sigma = \mathbf{K} + \sigma \mathbf{M}, \quad \text{with} \quad \sigma = \frac{2E}{\rho \ell^2}.
\]

\begin{verbatim}   
% Spectral shift
sigma = 2*E / (rho*l^2);

% Shifted stiffness matrix K_sigma
K_sigma = K + sigma * M;
\end{verbatim}

The symbolic matrices are evaluated numerically for specific material and geometric properties.

\begin{verbatim}   
% Numerical values
E_val = 210e9; 
A_val = 4e-4; 
rho_val = 7850; 
l_val = 0.5;

% Numerical evaluation
K_num = double(subs(K_sigma, {E, A, rho, l}, {E_val, A_val, rho_val, l_val}));
M_num = double(subs(M, {E, A, rho, l}, {E_val, A_val, rho_val, l_val}));
Lambda_exact_num = double(subs(Lambda_exact, {E, A, rho, l}, {E_val, A_val, 
        rho_val, l_val}));
P_exact_num = double(subs(P_exact, {E, A, rho, l}, {E_val, A_val, rho_val, l_val}));
\end{verbatim}

\textbf{3. Subspace iteration algorithm.}
The subspace iteration algorithm computes a subspace spanned by approximate eigenvectors. We perform three iterations of the subspace iteration algorithm starting from an initial guess $\mathbf P_0$. Each iteration involves:
\begin{enumerate}
   \item Solving $\mathbf{K}_\sigma\overline {\mathbf{P} }= \mathbf{M} \mathbf{P}$
   \item Projecting the problem onto the subspace: $\mathbf{K}_{proj}$ and $\mathbf{M}_{proj}$.
   \item Solving the reduced eigenproblem: $\mathbf \Lambda_{\sigma}$ and $\mathbf{Z}$,
   \item Reconstructing Ritz vectors: $\mathbf P$.
\end{enumerate}
Results are stored in cells P, Lambda, etc. The vectors are sorted to maintain order, and normalization is applied in the mass norm.
\begin{verbatim}   
    % Initial guess matrix
P0 = (1/sqrt(rho*A*l)) * [1 0; 0 1; 0 0];
P0_num = double(subs(P0, {E, A, rho, l}, {E_val, A_val, rho_val, l_val}));

% Definiton of matrices needed in the subspace iteration 
P = cell(1, 3);             % Store the approximated modal matrices
Pbar = cell(1, 3);          % Store the approximated non-orthogonal matrices
Lambda_sigma = cell(1, 3);  % Store the approximated offset spectral matrices
Lambda = cell(1, 3);        % Store the approximated spectral matrices
Z = cell(1, 3);             % Store the approximated projected modal matrices
K_proj = cell(1, 3);        % Store the projected stiffness matrices
M_proj = cell(1, 3);        % Store the projected mass matrices
% Start subspace iteration
for i = 1:3
    if i == 1
        Pbar{i} = K_num \ (M_num * P0_num);     
    else 
        Pbar{i} = K_num \ (M_num * P{i-1});
    end
   
    % Rayleigh matrices
    K_proj{i} = transpose(Pbar{i}) * K_num * Pbar{i};
    M_proj{i} = transpose(Pbar{i}) * M_num * Pbar{i};

    % Solve reduced eigenvalue problem
    [Z{i}, Lambda_sigma{i}] = eig(K_proj{i}, M_proj{i});
   
    % Update P (Ritz vectors)
    P{i} = Pbar{i} * Z{i};

    % Normalize mode shapes
    P{i} = P{i} / sqrt(transpose(P{i}) * M_num * P{i});  % M_num-norm normalization

    % Undo the spectral shift and consider only the real part
    Lambda{i} = real(diag(Lambda_sigma{i}) - double(subs(sigma, {E, rho, l}, 
            {E_val, rho_val, l_val})));
    Lambda_sigma{i} = real(diag(Lambda_sigma{i}));
    
    % Sort eigenvalues and apply permutation to eigenvectors
    [Lambda{i}, idx] = sort(Lambda{i});
    Lambda_sigma{i} = sort(Lambda_sigma{i});
    P{i} = real(P{i}(:, idx));
end
\end{verbatim}

\textbf{4) Comparison with exact solutions.} Exact eigenvalues and eigenvectors (up to scale) are known. We compare:
\begin{itemize}
   \item The eigenvalue approximation using relative error,
   \item The eigenvector approximation using the error in the $\mathbf M$-norm.
   \item The Ritz values converge to the exact ones and the approximate modes align with the analytical solutions.
\end{itemize}

\begin{verbatim}   
    %Initialize error arrays
rel_error_lambda = cell(1,3);
modal_error = cell(1,2);

% Loop over each iteration 
for j = 1:3
    % --- Relative error in eigenvalues ---
    rel_error_lambda{j} = abs(Lambda_sigma{j} - Lambda_exact_num) 
            ./ abs(Lambda_exact_num);

    % Loop over modes
    for m = 1:2
        % --- Modal error in M-norm ---
        p_diff = P{j}(:,m) - P_exact_num(:,m);
        modal_error{j}(m) = sqrt(p_diff' * M_num * p_diff);
    end
end    

for j = 1:3
    fprintf('\n======================================================\n');
    fprintf('\n--- Iteration %d ---\n', j);
    
    fprintf('Relative Error in Eigenvalues:\n');
    disp(rel_error_lambda{j});
  
    fprintf('Modal Error (M-norm):\n');
    disp(modal_error{j});
end
\end{verbatim}
The first eigenmode corresponds to a rigid-body motion, which has a zero natural frequency. The associated mode shape is a constant vector (e.g., \([1; 1; 1]\) or \([-1; -1; -1]\)), and it spans a one-dimensional eigenspace. Since eigenvectors are only defined up to a multiplicative constant—including sign—the subspace iteration algorithm may converge to either the positive or the negative of the exact mode shape.

When computing the modal error using the \( \mathbf{M} \)-norm (i.e., comparing the approximated and exact vectors directly), this ambiguity can result in misleading error values. Even though the vectors represent the same physical mode, their difference can have a nonzero norm if the signs differ. Specifically, the error between a normalized vector and its negative is twice the norm of the vector.

As a result, the modal error for the zero-frequency mode may appear to oscillate across iterations or fail to decrease monotonically, even though the convergence is mathematically correct. This behavior is not a flaw of the numerical method but a consequence of normalization and sign ambiguity. To mitigate this, one can minimize the error with respect to sign or explicitly align the direction of the approximated eigenvector with the reference one.

\end{document}
